% -*- coding: utf-8 -*-
% !TEX program = xelatex
\documentclass[plain,noanswer,medmath]{../../template/jnuexam}
\usepackage{../../template/kysx}

\begin{document}


\examtitle{1989}{3}


\loadproblems{1}{1989P1}%载入试卷一


\makepart{填空题}{本题满分15分，每小题3分}


\begin{problem}
$\lim_{x\to 0}\,x\cot 2x$= \fillin{}.
\end{problem}


\begin{problem}
$\int_0^{\pi}t\sin t\dt$=\fillin{}.
\end{problem}


\begin{problem}
曲线$y=\int_0^x (t-1)(t-2)\dt$在点$(0,0)$处的切线方程是\fillin{}.
\end{problem}


\begin{problem}
设$f(x)=x(x+1)(x+2)\cdots \cdots(x+n)$，则$f'(0)=$\fillin{}.
\end{problem}


% 使用试卷一，第一大题，第2小题
\useproblem{1}{1}{2}


\begin{problem}
设$f(x)=\begin{cases}
  a+bx^2,x\le 0 \\
  \frac{\sin bx}{x},x>0
\end{cases}$在$x=0$处连续，则常数$a$与$b$应满足的关系是\fillin{}.
\end{problem}


\begin{problem}
设$\tan y=x+y$，则$\dy=$\fillin{}.
\end{problem}


\makepart{计算题}{本题满分20分，每小题4分}


\begin{problem}
已知$y=\arcsin\e^{-\sqrt{x}}$，求$y'$．
\end{problem}


\begin{problem}
求$\int \frac{\dx}{x\ln^2 x}$．
\end{problem}


\begin{problem}
求$\lim_{x\to 0}\,(2\sin x+\cos x)^{\frac{1}{x}}$．
\end{problem}


\begin{problem}
已知$\begin{cases}
  x=\ln(1+t^2), \\
  y=\arctan t, \\
\end{cases}$求$\frac{\dy}{\dx}$及$\frac{\d^2y}{\dx^2}$．
\end{problem}


\begin{problem}
已知$f(2)=\frac{1}{2}$,$f'(2)=0$及$\int_0^2f(x)\dx=1$，求$\int_0^1x^2f''(2x)\dx$．
\end{problem}


\makepart{选择题}{本题满分18分，每小题3分}


\useproblem{1}{2}{1}


\begin{problem}
若$3a^2-5b<0$，则方程$x^5+2ax^3+3bx+4c=0$ \pickout{}
\begin{abcd}
\item 无实根．
\item 有唯一实根．
\item 有三个不同实根．
\item 有五个不同实根．
\end{abcd}
\end{problem}


\begin{problem}
曲线$y=\cos x$（$-\frac{\pi}{2}\le x\le \frac{\pi}{2}$）与$x$轴所围成的图形，
绕$x$轴旋转一周所成的旋转体的体积为 \pickout{}
\begin{abcd}
\item $\frac{\pi}{2}$．
\item $\pi$．
\item $\frac{\pi^2}{2}$．
\item $\pi^2$．
\end{abcd}
\end{problem}


\begin{problem}
设两函数$f(x)$及$g(x)$都在$x=a$处取得极大值，
则函数$F(x)=f(x)g(x)$在$x=a$处 \pickout{}
\begin{abcd}
\item 必取极大值．
\item 必取极小值．
\item 不可能取极值．
\item 是否取极值不能确定．
\end{abcd}
\end{problem}


\begin{problem}
微分方程$y''-y=\e^x+1$的一个特解应具有形式（式中$a,b$为常数） \pickout{}
\begin{abcd}
\item $a\e^x+b$．
\item $ax\e^x+b$．
\item $a\e^x+bx$．
\item $ax\e^x+bx$．
\end{abcd}
\end{problem}


\begin{problem}
设$f(x)$在$x=a$的某个邻域内有定义，则$f(x)$在$x=a$可导的一个充分条件是 \pickout{}
\begin{abcd}
\item $\lim_{h\to +\infty}\,h[f(a+\frac{1}{h})-f(a)]$存在．
\item $\lim_{h\to 0}\,\frac{f(a+h)-f(a+h)}{h}$存在．
\item $\lim_{h\to 0}\,\frac{f(a+h)-f(a-h)}{2h}$存在．
\item $\lim_{h\to 0}\,\frac{f(a)-f(a-h)}{h}$存在．
\end{abcd}
\end{problem}


\makepart{}{本题满分6分}

\begin{problem}
求微分方程$xy'+(1-x)y=\e^{2x}$($0<x<+\infty$)满足$y(1)=0$的解．
\end{problem}


\makepart{}{本题满分7分}

\useproblem{1}{5}{0}%【同试卷一第五题】


\makepart{}{本题满分7分}

\useproblem{1}{6}{0}%【同试卷一第六题】


\makepart{}{本题满分11分}

\begin{problem}
对函数$y=\frac{x+1}{x^2}$，填写下表： \par
\begin{tabularx}{\linewidth}{|>{\qquad}X|>{\qquad}X|}
\hline
  单调减少区间 &  \\
\hline
  单调增加区间 &  \\
\hline
  极值点 &  \\
\hline
  极值 &  \\
\hline
  凹$\left(\bigcup\right)$区间 &  \\
\hline
  凸$\left(\bigcap\right)$区间 &  \\
\hline
  拐点 &  \\
\hline
  渐近线 &  \\
\hline
\end{tabularx}
\end{problem}


\makepart{}{本题满分10分}

\begin{problem}
设抛物线$y=ax^2+bx+c$过原点，当$0\le x\le 1$时$y\ge 0$，
又已知该抛物线与$x$轴及直线$x=1$所围图形的面积为$\frac{1}{3}$，
试确定$a,b,c$使此图形绕$x$旋转一周而成的旋转体的体积$V$最小．
\end{problem}


\end{document}
